% Management in Global Markets – Midterm Study Guide — EU business casual
% Compile with: pdflatex study-guide.tex

\documentclass[8pt,landscape,a4paper]{article}
\usepackage[utf8]{inputenc}
\usepackage[T1]{fontenc}
\usepackage[left=0.8cm,right=2.6cm,top=0.8cm,bottom=0.8cm,marginparwidth=2cm,marginparsep=0.3em]{geometry}
\usepackage{multicol}
\usepackage{marginnote}
\usepackage{enumitem}
\setlist{nosep,leftmargin=*}
\usepackage{booktabs}
\usepackage{parskip}
\usepackage{xcolor}
\usepackage{fancyhdr}
\definecolor{eublue}{RGB}{0,51,153}
\definecolor{charcoal}{RGB}{51,51,51}
\definecolor{softgrey}{RGB}{240,240,240}
\definecolor{notegrey}{RGB}{90,90,90}
\newcommand{\handnote}[1]{\marginnote{\raggedright\footnotesize\color{notegrey}\textit{#1}\vspace{0.4em}}}
\newcommand{\hlpink}[1]{\textcolor{charcoal}{\textbf{#1}}}
\newcommand{\hlpeach}[1]{\textcolor{eublue}{\textbf{#1}}}
\setlength{\parskip}{0.2em}
\setlength{\columnsep}{0.5em}
\pagestyle{fancy}
\fancyhf{}
\fancyfoot[L]{\raisebox{0.8em}{\small\color{notegrey} Mgmt Global}}
\fancyfoot[R]{\raisebox{0.8em}{\small\color{notegrey}\thepage}}
\fancypagestyle{plain}{\fancyhf{}\fancyfoot[L]{\raisebox{0.8em}{\small\color{notegrey} Mgmt Global}}\fancyfoot[R]{\raisebox{0.8em}{\small\color{notegrey}\thepage}}}
\raggedright

\begin{document}
\begin{center}
\vspace{0.2em}
{\normalsize\bfseries\color{eublue} Management in Global Markets -- Midterm Study Guide}\\[0.2em]
\footnotesize\textcolor{notegrey}{Use for quick recall; in exam: define terms, give one example, tie to internationalization.}\\[0.35em]
\noindent\colorbox{softgrey}{\parbox{0.92\textwidth}{\centering\footnotesize\color{charcoal}
\textbf{Key:} \textbf{Definitions} \quad \textcolor{eublue}{\textbf{Frameworks \& tips}} \quad \textit{Margin notes}
\vspace{0.25em}}}
\end{center}
\vspace{-0.4em}

\begin{multicols}{2}

\section*{\color{eublue}Unit 0 -- Course \& Big Ideas}\handnote{Location = masterpiece.}
\begin{itemize}[leftmargin=*,nosep]
\item \textbf{International business} = Operating across borders. Goods, services, capital, people move.
\item Key question: Why go abroad? Market size, resources, efficiency, learning.
\item Reality: Most companies fail. Wrong location, weak CA, poor execution.
\item \textbf{\hlpink{``Location is the masterpiece''}} = Choosing right location $>$ skills, timing, hard work.
\end{itemize}

\section*{\color{eublue}Unit 1 -- Globalization vs.\ Slowbalization}\handnote{Winners vs.\ losers.}
\begin{itemize}[leftmargin=*,nosep]
\item \textbf{\hlpink{Globalization}} = Free movement. High production + export. Interconnected economies.
\item \textbf{\hlpink{Slowbalization}} = Slowing. Protectionist. Localization.
\item Winners: China (biggest exporter), India, Indonesia. USA = biggest importer.
\item \textbf{\hlpink{Economic complexity}} = Raw $\to$ complex products (tech, R\&D). USA, UK, Netherlands strong.
\item Demographics: Europe 60\% over 60 by 2032. Africa growing. Young = production.
\item Protectionism: Tariffs. G20: $\sim$50\% global economy. Spain in G20.
\item \textbf{Localization} = Adapt for local markets. Before internationalization.
\end{itemize}\handnote{China = exporter. USA = importer.}

\section*{\color{eublue}Unit 2 -- Supply Chain \& Products}\handnote{Quality vs.\ cost.}
\begin{itemize}[leftmargin=*,nosep]
\item \textbf{\hlpink{Supply chain}} = Parts move before product made. Global = many countries.
\item Japan = quality, higher price. China = cost, volume, quality varies.
\item Products: shoes (Vietnam), shirt (India), watch (China). Rarely 100\% one country.
\item Aircraft: parts from multiple countries, assembled in one place (e.g.\ Chicago).
\item \textbf{Risks}: Cultural, legal, logistical, financial, communication. ``You cannot escape language.''
\item Market attractiveness: GDP growth, population growth, political risk.
\item \textbf{Go where competitors are} (Option B). Validation, infrastructure, learn. Not Option C (alone).\handnote{Competitors = market works!}
\item Culture shock: honeymoon $\to$ hard $\to$ adapt. Inevitable.
\end{itemize}

\section*{\color{eublue}Unit 3 -- Hofstede's Dimensions}\handnote{Culture = cannot escape.}
\begin{itemize}[leftmargin=*,nosep]
\item \textbf{Culture} = Shared beliefs, values, norms. Distinguishes groups.
\item \textbf{Power distance}: High (Spain, Mexico) vs.\ low (USA, N.\ Europe).
\item \textbf{Individualism vs.\ collectivism}: Rich = individualistic; poorer = collective.
\item \textbf{Uncertainty avoidance}: Comfort with ambiguity.
\item \textbf{Long-term vs.\ short-term}: Security $\to$ plan; uncertainty $\to$ survive today.
\item Key: Must adapt. Cannot escape.
\end{itemize}\handnote{Five dimensions. Know at least 2--3.}

\section*{\color{eublue}Unit 4 -- Four Phases \& Location}\handnote{Phase 3 = most dangerous.}
\begin{itemize}[leftmargin=*,nosep]
\item \textbf{Phase 1 -- Conceive}: Plan, research. 2--3 years.
\item \textbf{Phase 2 -- Execute}: Move, plant. 2 years. Highest investment. \textbf{Critical}---wrong place = fail.
\item \textbf{Phase 3 -- Establish}: Running, making money. \textbf{Most dangerous}---most fail here.
\item \textbf{Phase 4 -- Expand}: Success. Scale.
\item Total: 4--5 years.
\item \textbf{Artificial risk} (laws, rules): can manage. \textbf{Natural risk} (language, culture): cannot escape; adapt.
\item Why fail: wrong location, poor communication, no backup, left alone.
\end{itemize}

\section*{\color{eublue}Unit 5 -- Strategy Levels \& Int'l Logic}\handnote{Corporate, competitive, functional.}
\begin{itemize}[leftmargin=*,nosep]
\item \textbf{Corporate}: Which businesses? (Cars vs.\ trucks).
\item \textbf{Competitive}: How to compete? (Price, quality, niche).
\item \textbf{Functional}: HR, marketing, finance, ops, R\&D.
\item Keep simple. 15 years, not 15 days. ``Slow is fast.''
\item \textbf{\hlpink{SME}} = $<$250 employees, $<$50M. Spain: 99\%. Few internationalize.
\item \textbf{Location choice}: Critical. Bad env = fail regardless of internal strengths.
\item \textbf{Resources} = what you have. \textbf{Capabilities} = what you do.\handnote{Have vs.\ do.}
\item \textbf{Competitive advantage} = better than rivals. \textbf{Absolute} = almost no one has. Rare.
\item Before abroad: strong CA? If ``just like'' others in mature market = very risky.
\item Entry: acquisition, JV, alliances (licensing, distribution). Partner = local knowledge.
\end{itemize}\handnote{Before abroad: strong CA? If not, stay home.}

\subsection*{Resources vs.\ capabilities}
\begin{center}
\footnotesize
\begin{tabular}{@{}lll@{}}
\toprule
\textbf{Aspect} & \textbf{Resources} & \textbf{Capabilities} \\
\midrule
What & What you \textit{have} & What you \textit{do} \\
Examples & Assets, people, brand & Processes, routines \\
\bottomrule
\end{tabular}
\end{center}

\section*{\color{eublue}Unit 6 -- Chile Case}\handnote{Tech transfer risk.}
\begin{itemize}[leftmargin=*,nosep]
\item French company $\to$ Chile. Capital, people, tech. Shared with local partners.
\item Failed: lost money, jobs, tech. Partners used tech independently. No protection. Communication broke.
\item Lesson: careful with core tech. IP, contracts, partners. Culture, trust.
\end{itemize}

\section*{\color{eublue}Unit 7 -- Export vs.\ Int'l \& FDI}\handnote{Product vs.\ company.}
\begin{itemize}[leftmargin=*,nosep]
\item \textbf{\hlpink{Export}} = Produce home. Ship abroad. Low risk. Limited commitment.
\item \textbf{\hlpink{Internationalization}} = Plants, subsidiaries, HQ abroad. Needs FDI, people.
\item \textbf{\hlpink{FDI}} = Long-term commitment. Creates need for international HRM.
\item \textbf{Domestic}: all in one country. \textbf{International}: several countries, few key locations.
\item \textbf{Transnational}: global brands, adapt locally (cosmetics, food, fashion).
\end{itemize}\handnote{Export = first step. FDI = next.}

\section*{\color{eublue}Unit 8 -- International HRM}\handnote{People matter.}
\begin{itemize}[leftmargin=*,nosep]
\item \textbf{International HRM} = Multiple countries, cultures. Cross-border mobility. Coordinate policies.
\item First questions: Which functions? How many people? Have them or hire? Support families?
\item \textbf{Who goes?} Younger, flexible. Internal preferred for sensitive (trust). External = riskier.
\item \textbf{Talent} = Skills + availability to move. Some stage may not want to move.
\item \textbf{Motivation}: Salary + ambition, recognition, quality of life. As important as skills.
\item Turnover: young leave. Negotiate. Know market value.
\item Trust: prefer known people. ``Keys to the kingdom'' to outsiders = risky.
\end{itemize}

\section*{\color{eublue}Unit 9 -- Cash Generation}\handnote{Cash from customers.}
\begin{itemize}[leftmargin=*,nosep]
\item Cash comes from customers. No sell = no cash. Business = profit.
\item Two options: minimize expenses OR push sales.
\item Repeat business = most important. Communication. Deliver on promises.
\end{itemize}\handnote{Two options only.}

\section*{\color{eublue}Quick revision -- Can you explain\ldots?}\handnote{Run through before exam.}
\begin{itemize}[leftmargin=*,nosep]
\item Global vs.\ slowbalization. Economic complexity.

\item Supply chain. Japan vs.\ China. Products from many countries.
\item Risks. Why go where competitors are. Culture shock.
\item Hofstede: power distance, individualism, long-term orientation.
\item Four phases. Which most dangerous? Artificial vs.\ natural risk.
\item Strategy levels. SME. Location. Resources vs.\ capabilities.
\item CA vs.\ absolute. Chile case. Export vs.\ int'l vs.\ FDI.
\item Domestic vs.\ international vs.\ transnational.
\item International HRM. Who goes. Motivation. Trust.
\item Cash generation.
\end{itemize}\handnote{If you can explain all these, you're ready.}

\section*{\hlpeach{Common question angles}}\handnote{Memorise these.}
\begin{itemize}[leftmargin=*,nosep]
\item \textbf{``Why might int'l fail?''} Wrong location, weak CA, tech transfer (Chile), culture, Phase 2--3 execution.
\item \textbf{``Strategy levels?''} Corporate / competitive / functional.
\item \textbf{``Export vs.\ int'l?''} Product moves vs.\ company moves operations.
\item \textbf{``Success abroad?''} Strong CA, location, people, careful tech, partners, go where competitors are.
\item \textbf{``Why competitors?''} Validation, infrastructure, learn, differentiate.
\item \textbf{``Apply four phases.''} Conceive $\to$ Execute (critical) $\to$ Establish (dangerous) $\to$ Expand.
\end{itemize}

\end{multicols}

\newpage
\vspace*{0.3em}
\begin{center}
\textbf{\color{eublue}Quick recap -- last page}\\[0.3em]
\small\textcolor{notegrey}{(Section summaries, not margin notes)}
\end{center}
\vspace{0.4em}
\noindent\color{notegrey}
\setlength{\parskip}{0.3em}
\begin{multicols}{2}

\textbf{Unit 0 -- Big Ideas}\\
Location = masterpiece. Most fail.

\textbf{Unit 1 -- Globalization}\\
Global vs.\ slowbalization. Economic complexity. Winners vs.\ losers.

\textbf{Unit 2 -- Supply Chain}\\
Quality vs.\ cost. Japan vs.\ China. Go where competitors are. Culture shock.

\textbf{Unit 3 -- Hofstede}\\
Power distance. Individualism. Long-term. Cannot escape culture.

\textbf{Unit 4 -- Four Phases}\\
Conceive. Execute (critical). Establish (most dangerous). Expand. Artificial vs.\ natural risk.

\textbf{Unit 5 -- Strategy}\\
Corporate. Competitive. Functional. SME. Location. Resources vs.\ capabilities. CA vs.\ absolute.

\textbf{Unit 6 -- Chile}\\
Tech transfer risk. No protection. Communication broke.

\textbf{Unit 7 -- Export vs.\ Int'l}\\
Product vs.\ company. FDI. Domestic vs.\ international vs.\ transnational.

\textbf{Unit 8 -- HRM}\\
Who goes. Internal vs.\ external. Motivation. Trust. Turnover.

\textbf{Unit 9 -- Cash}\\
From customers. Two options. Repeat business.

\textbf{Common angles}\\
Why fail. Strategy levels. Export vs.\ int'l. Success. Four phases.

\end{multicols}
\vfill
\normalsize\normalfont\color{black}
\end{document}
